\documentclass{report}
\usepackage{amsmath,amssymb}
\setlength{\parindent}{0mm}
\setlength{\parskip}{1em}
\begin{document}
\begin{center}
\rule{6in}{1pt} \
{\large Christian Glusa \\
{\bf Resilience Analysis for Multigrid Methods}}

Division of Applied Mathematics \\ 182 George Street \\ Brown University \\ Providence \\ RI 02912
\\
{\tt christian\_glusa@brown.edu}\\
Mark Ainsworth\end{center}

With the the advent of exascale computing expected within the next few
years, the number of components in a system will continue to grow. The
error rate per individual component is unlikely to improve, however,
meaning that future high performance computing will be faced with faults
occurring at significantly higher rates than present day installations.
Therefore, the resilience properties of numerical methods will become
important factors in both the choice of algorithm and in its analysis.

In this talk we present a framework for the analysis of linear iterative
methods in a fault-prone environment. The effects of random node failures
are taken into account through a probabilistic model involving random
diagonal matrices. Using this model, we analyze the behavior of two- and
multigrid methods under random node failures. Our results show that while
standard multigrid is not resilient, protecting the prolongation leads to
a fault-resilient variant. Both analytic convergence estimates for these
methods and simulation results will be discussed.

This is joint work with Mark Ainsworth.


\end{document}
